% ======================================================================
%  SIMC 2026 --- PRELIMINARY WORKING REPORT  \textperiodcentered{}  Questions 1 & 2 only
%  Endeavour Track  \textperiodcentered{}  INTERNAL DRAFT --- full proofs, all details.
%
%  Source of all mathematical content:
%     notes/Note 26 May 2026.pdf      (Q1, parts a--b)
%     notes/Note 26 May 2026 (2).pdf  (Q1, parts a--e, g)
%     notes/task 2.pdf                (Q2, parts a--e)
%
%  Every deviation from those notes (mistake fixed, gap filled, extra
%  added, elaboration) is logged in  extras.md  in this folder.
%
%  Team:  K. Vishnu Kumar \textperiodcentered{} Adithya Kesan Jayakanth \textperiodcentered{} Adithya Anand
%  Date:  26 May 2026
% ======================================================================

\documentclass[12pt,a4paper]{article}

\usepackage[a4paper, margin=2.0cm, headheight=14pt]{geometry}
\usepackage{mathptmx}                 % Times for text + math
\usepackage[T1]{fontenc}
\usepackage{amsmath,amssymb,amsthm,bm}
\usepackage{enumitem}
\usepackage{xcolor}
\usepackage{fancyhdr}
\usepackage[hidelinks]{hyperref}

\linespread{1.10}
\setlength{\parskip}{5pt}
\setlength{\parindent}{0pt}

% --- header / footer --------------------------------------------------
\pagestyle{fancy}
\fancyhf{}
\fancyhead[L]{\small SIMC 2026 \textperiodcentered{} Preliminary Working Report (Q1 \& Q2)}
\fancyhead[R]{\small \thepage}
\fancyfoot[C]{\small Vishnu \textperiodcentered\ Adithya K.J. \textperiodcentered\ Adithya A. \quad\textbf{INTERNAL DRAFT --- NOT FOR SUBMISSION}}
\renewcommand{\headrulewidth}{0.4pt}

% --- review flags (mirrored in extras.md) -----------------------------
\newcommand{\miss}[1]{{\textbf{\textcolor{red}{[MISS: #1]}}}}
\newcommand{\fix}[1]{{\textbf{\textcolor{orange!85!black}{[CORRECTED: #1]}}}}
\newcommand{\extra}[1]{{\textbf{\textcolor{blue!80!black}{[EXTRA: #1]}}}}

% --- proof environments -----------------------------------------------
\theoremstyle{definition}
\newtheorem*{claim}{Claim}

% --- notation shortcuts -----------------------------------------------
\newcommand{\R}{\mathbb{R}}
\newcommand{\Mx}{\mathbf{M}}
\newcommand{\Xb}{\mathbf{X}}
\newcommand{\vv}{\mathbf{v}}
\newcommand{\uu}{\mathbf{u}}
\newcommand{\xx}{\mathbf{x}}
\newcommand{\yy}{\mathbf{y}}
\newcommand{\pp}{\mathbf{p}}
\newcommand{\rr}{\mathbf{r}}
\newcommand{\zero}{\mathbf{0}}
\newcommand{\norm}[1]{\lVert #1 \rVert}
\newcommand{\T}{^{\!\top}}

\title{\vspace{-1.6cm}\Huge\bfseries A Linear Algebra Guide to the\\Strange Life of High-Dimensional Data\\[0.4em]
\Large Preliminary Working Report --- Questions 1 \& 2}
\author{K.\ Vishnu Kumar \and Adithya Kesan Jayakanth \and Adithya Anand}
\date{26 May 2026\\[0.3em]\small\itshape Internal draft: full proofs and all details.
Questions 1 and 2 only; to be consolidated with Question 3 into the 10-page submission later.}

\begin{document}
\maketitle
\thispagestyle{fancy}

% =====================================================================
\section*{Scope and status of this draft}
\addcontentsline{toc}{section}{Scope and status}

This preliminary report contains the \emph{complete} worked solutions to
\textbf{Question 1} (Linear Algebra of Markov Matrices) and \textbf{Question 2}
(The Power Method and Singular Value Decomposition), with every proof written
out in full. It is a working document, not the final submission: the
generic report scaffolding (literature, data provenance, sensitivity studies)
is omitted because Questions 1 and 2 are purely analytical, and Question 3
(the data analysis on \texttt{mystery\_matrix1.npy}) will be merged in later
before the work is compressed to the 10-page cap.

Coloured tags mark every place this document departs from the team's
handwritten notes. Each tag has a matching entry in \texttt{extras.md}:
\begin{itemize}[leftmargin=1.4em, itemsep=1pt]
  \item \miss{\dots} --- a part of the problem not yet attempted in the notes.
  \item \fix{\dots} --- a mathematical error in the notes, corrected here.
  \item \extra{\dots} --- content added beyond the notes (an elaboration or a missing explanation).
\end{itemize}

% =====================================================================
\section*{Notation and conventions}
\addcontentsline{toc}{section}{Notation and conventions}

\begin{tabular}{@{}p{0.20\linewidth}p{0.74\linewidth}@{}}
$\Mx,\ \Xb$ & matrices (bold uppercase). \\
$\vv,\uu,\xx,\pp$ & column vectors (bold lowercase). \\
$\Xb\T$ & transpose of $\Xb$, $(\Xb\T)_{ij}=\Xb_{ji}$. \\
$\lambda_i,\ \vv_i$ & eigenvalue / eigenvector pair, ordered $|\lambda_1|\ge|\lambda_2|\ge\cdots$. \\
$\sigma_i$ & singular value, $\sigma_i=\sqrt{\lambda_i}$. \\
$\norm{\xx}$ & Euclidean norm, $\norm{\xx}=\sqrt{\xx\T\xx}$. \\
\end{tabular}

A set of vectors is \emph{orthonormal} if $\vv_i\T\vv_j=0$ for $i\ne j$ and
$\vv_i\T\vv_i=1$. We use the \textbf{spectral theorem} as a given: every real
symmetric matrix has real eigenvalues and an orthonormal basis of real
eigenvectors.

% =====================================================================
% =====================================================================
\section{Linear Algebra of Markov Matrices}

\textbf{Setup.} Two organism types have fractional populations
$p_1(t),p_2(t)$, evolving by $\pp(t+1)=\Mx\,\pp(t)$ with
\[
\pp(t)=\begin{pmatrix}p_1(t)\\ p_2(t)\end{pmatrix},
\qquad
\Mx=\begin{pmatrix}1-\varepsilon & 0\\[2pt] \varepsilon & 1\end{pmatrix},
\qquad \varepsilon\in(0,1)\ \text{fixed.}
\]
Componentwise this is
\begin{equation}\label{eq:update}
p_1(t+1)=(1-\varepsilon)\,p_1(t),\qquad
p_2(t+1)=\varepsilon\,p_1(t)+p_2(t).
\end{equation}
Each column of $\Mx$ sums to $1$ (it is \emph{column-stochastic}); a fraction
$\varepsilon$ of type-1 converts to type-2 each generation.

% ---------------------------------------------------------------------
\subsection{(a) Conservation}

\begin{claim}
$p_1(t)+p_2(t)$ is constant in $t$.
\end{claim}
\begin{proof}
Let $S(t)=p_1(t)+p_2(t)$. Adding the two equations in \eqref{eq:update},
\[
S(t+1)=p_1(t+1)+p_2(t+1)=(1-\varepsilon)p_1(t)+\bigl(\varepsilon\,p_1(t)+p_2(t)\bigr)
       = p_1(t)+p_2(t)=S(t).
\]
So $S(t+1)=S(t)$ for every $t$. By induction $S(t)=S(0)=p_1(0)+p_2(0)$, a
constant independent of $t$. \end{proof}

\textit{Equivalent one-line argument.} Because every column of $\Mx$ sums to
$1$, the row vector $\mathbf{1}\T=(1,1)$ satisfies $\mathbf{1}\T\Mx=\mathbf{1}\T$,
hence $\mathbf{1}\T\pp(t+1)=\mathbf{1}\T\Mx\,\pp(t)=\mathbf{1}\T\pp(t)$, i.e.\
$S(t+1)=S(t)$.

% ---------------------------------------------------------------------
\subsection{(b) Eigenspectrum}

The eigenvalues solve $\det(\Mx-\lambda \mathbf I)=0$:
\[
\det\!\begin{pmatrix}1-\varepsilon-\lambda & 0\\[2pt] \varepsilon & 1-\lambda\end{pmatrix}
=(1-\varepsilon-\lambda)(1-\lambda)=0
\quad\Longrightarrow\quad
\boxed{\lambda_1=1,\qquad \lambda_2=1-\varepsilon.}
\]
(The matrix is lower-triangular, so the eigenvalues are simply its diagonal
entries.) The eigenvectors come from $(\Mx-\lambda_i\mathbf I)\vv_i=\zero$.

\textit{For $\lambda_1=1$:}\quad
$\begin{pmatrix}-\varepsilon & 0\\ \varepsilon & 0\end{pmatrix}\!\begin{pmatrix}x\\ y\end{pmatrix}=\zero
\Rightarrow -\varepsilon x=0\Rightarrow x=0$, $y$ free. Take
$\displaystyle \vv_1=\begin{pmatrix}0\\1\end{pmatrix}.$

\textit{For $\lambda_2=1-\varepsilon$:}\quad
$\begin{pmatrix}0 & 0\\ \varepsilon & \varepsilon\end{pmatrix}\!\begin{pmatrix}x\\ y\end{pmatrix}=\zero
\Rightarrow \varepsilon(x+y)=0\Rightarrow y=-x$. Take
$\displaystyle \vv_2=\begin{pmatrix}1\\-1\end{pmatrix}.$

% ---------------------------------------------------------------------
\subsection{(c) Linear combination}

Assume $p_1(0)=\tfrac12$. Write
$\pp(0)=c_1\vv_1+c_2\vv_2=c_1\!\begin{pmatrix}0\\1\end{pmatrix}+c_2\!\begin{pmatrix}1\\-1\end{pmatrix}
=\begin{pmatrix}c_2\\ c_1-c_2\end{pmatrix}.$
Matching components,
\[
c_2=p_1(0)=\tfrac12,\qquad c_1-c_2=p_2(0)\ \Rightarrow\ c_1=p_2(0)+\tfrac12 .
\]
\fix{notes assumed $p_2(0)=\tfrac12$ without stating it.} If, as is natural for
fractional populations, the totals are normalised so that
$p_1(0)+p_2(0)=1$, then $p_2(0)=\tfrac12$ and $c_1=1$, giving
\[
\boxed{\ \pp(0)=\vv_1+\tfrac12\,\vv_2.\ }
\]
More generally $c_1=p_2(0)+\tfrac12$ and $c_2=\tfrac12$; note
$c_1=p_1(0)+p_2(0)=S(0)$, the conserved total from part (a).

% ---------------------------------------------------------------------
\subsection{(d) Eigenvectors make time evolution easy}

With $\pp(0)=c_1\vv_1+c_2\vv_2$, apply $\Mx$ repeatedly. Since
$\Mx\vv_i=\lambda_i\vv_i$, we have $\Mx^{t}\vv_i=\lambda_i^{\,t}\vv_i$, so
\[
\pp(t)=\Mx^{t}\pp(0)=c_1\lambda_1^{\,t}\vv_1+c_2\lambda_2^{\,t}\vv_2
\quad\Longrightarrow\quad
\boxed{\ \pp(t)=c_1\,\vv_1+c_2\,(1-\varepsilon)^{t}\,\vv_2.\ }
\]
The eigenbasis turns repeated matrix multiplication into scalar powers: each
mode evolves independently by its own factor $\lambda_i^{\,t}$.

% ---------------------------------------------------------------------
\subsection{(e) Long-time behaviour}

\textbf{(i)}~Because $\varepsilon\in(0,1)$ we have $0<1-\varepsilon<1$, so
$(1-\varepsilon)^{t}\to0$ as $t\to\infty$. Hence
\[
\pp(t)\ \xrightarrow{\ t\to\infty\ }\ c_1\vv_1,
\]
the eigenvector with the larger eigenvalue $\lambda_1=1>\lambda_2$.

\textbf{(ii)}~The structure of $\Mx$ explains this. The second column
$(0,1)\T$ means type-2 is \emph{absorbing}: nothing leaves it. Type-1 leaks
into type-2 at rate $\varepsilon$ every step and never receives anything back.
So all mass drains into type-2. Quantitatively the limit is
\[
c_1\vv_1=\begin{pmatrix}0\\ c_1\end{pmatrix}
=\begin{pmatrix}0\\ p_1(0)+p_2(0)\end{pmatrix},
\]
i.e.\ the entire conserved population $S(0)$ ends up as type-2. The eigenvalue
$\lambda_1=1$ carries the conserved steady state, while $\lambda_2=1-\varepsilon<1$
governs the decaying transient.

% ---------------------------------------------------------------------
\subsection{(f) Validation}

\miss{not attempted in the notes.} The closed form
$\pp(t)=c_1\vv_1+c_2(1-\varepsilon)^t\vv_2$ should be checked against a direct
numerical iteration of $\pp(t+1)=\Mx\pp(t)$ from $\pp(0)=(\tfrac12,\tfrac12)$,
for a sample $\varepsilon$ (e.g.\ $0.1$), confirming agreement at every $t$ and
the limit $\pp(\infty)=(0,1)$. \emph{To be added in} \texttt{src/task1.py}
\emph{with a convergence plot; no numbers are invented here.}

% ---------------------------------------------------------------------
\subsection{(g) Time-varying conversion rate}

Now $\varepsilon(t)=\varepsilon_0\,e^{-t/t_0}$ with $\varepsilon_0,t_0>0$, so the
update matrix changes each step:
$\Mx(t)=\bigl(\begin{smallmatrix}1-\varepsilon(t) & 0\\ \varepsilon(t) & 1\end{smallmatrix}\bigr)$.

\fix{the notes wrote $\pp(t)=1^{t}c_1+\bigl(1-\varepsilon(t)\bigr)^{t}c_2$, which is
not correct for a time-varying rate.} The key observation is that the
eigenvectors $\vv_1=(0,1)\T$ and $\vv_2=(1,-1)\T$ \emph{do not depend on
$\varepsilon$}, so they are shared by \emph{every} matrix $\Mx(t)$:
\[
\Mx(t)\,\vv_1=1\cdot\vv_1,\qquad \Mx(t)\,\vv_2=\bigl(1-\varepsilon(t)\bigr)\vv_2 .
\]
The state is the ordered product $\pp(t)=\Mx(t-1)\Mx(t-2)\cdots\Mx(0)\,\pp(0)$.
Acting on $\pp(0)=c_1\vv_1+c_2\vv_2$, the $\vv_1$ mode is multiplied by $1$ at
every step, while the $\vv_2$ mode picks up one factor $\bigl(1-\varepsilon(s)\bigr)$
per step. Therefore
\[
\boxed{\ \pp(t)=c_1\,\vv_1+c_2\Bigl[\textstyle\prod_{s=0}^{t-1}\bigl(1-\varepsilon_0 e^{-s/t_0}\bigr)\Bigr]\vv_2.\ }
\]
So the constant-rate factor $(1-\varepsilon)^{t}$ of part (d) is replaced by the
\textbf{product} $\prod_{s=0}^{t-1}\bigl(1-\varepsilon(s)\bigr)$, not by
$\bigl(1-\varepsilon(t)\bigr)^{t}$.

\extra{qualitative consequence (beyond the notes).} Since
$\sum_{s\ge0}\varepsilon_0 e^{-s/t_0}<\infty$, the infinite product converges to a
\emph{strictly positive} limit $\Pi_\infty=\prod_{s=0}^{\infty}\bigl(1-\varepsilon_0 e^{-s/t_0}\bigr)$.
Unlike the constant-$\varepsilon$ case, the $\vv_2$ transient does \emph{not}
decay to zero: the conversion ``switches off'' as $\varepsilon(t)\to0$, freezing a
residual amount $c_2\Pi_\infty$ in the $\vv_2$ direction. The system therefore
settles at $\pp(\infty)=c_1\vv_1+c_2\Pi_\infty\vv_2$, i.e.\ type-1 is
\emph{not} fully depleted.

% ---------------------------------------------------------------------
\subsection{(h) Validate again}

\miss{not attempted in the notes.} The corrected product formula for part (g)
should be validated the same way as (f): iterate $\pp(t+1)=\Mx(t)\pp(t)$ with
$\varepsilon(t)=\varepsilon_0 e^{-t/t_0}$ and compare against
$c_1\vv_1+c_2\prod_{s<t}(1-\varepsilon(s))\,\vv_2$, checking the residual
$c_2\Pi_\infty$ in the limit. \emph{To be added in} \texttt{src/task1.py}.

% =====================================================================
% =====================================================================
\section{The Power Method and Singular Value Decomposition}

\textbf{Setup.} $\Xb$ is a symmetric $n\times n$ matrix with eigenvalues
$|\lambda_1|>|\lambda_2|\ge\cdots\ge|\lambda_n|$ and an orthonormal eigenbasis
$\vv_1,\dots,\vv_n$. The Power Method starts from a unit vector $\xx_0$ and
iterates
\begin{equation}\label{eq:power}
\xx_{k+1}=\frac{\Xb\,\xx_k}{\norm{\Xb\,\xx_k}},\qquad k=0,1,2,\dots
\end{equation}

% ---------------------------------------------------------------------
\subsection{(a) The Power Method converges}

\begin{claim}
If $\xx_0$ has a nonzero component along $\vv_1$, then $\xx_k\to\pm\vv_1$.
\end{claim}
\begin{proof}
Normalisation in \eqref{eq:power} only rescales, so $\xx_k$ has the same
\emph{direction} as the unnormalised iterate
$\yy_k:=\Xb^{k}\xx_0$. Expand $\xx_0$ in the orthonormal eigenbasis,
\[
\xx_0=\sum_{i=1}^{n} a_i\vv_i,\qquad a_i=\xx_0\T\vv_i,\qquad a_1\ne0 .
\]
Since $\Xb^{k}\vv_i=\lambda_i^{k}\vv_i$,
\[
\yy_k=\Xb^{k}\xx_0=\sum_{i=1}^{n} a_i\lambda_i^{k}\vv_i
      =\underbrace{a_1\lambda_1^{k}\vv_1}_{\text{parallel to }\vv_1}
      +\underbrace{\sum_{i=2}^{n} a_i\lambda_i^{k}\vv_i}_{=\ \rr_k\ \text{(residual)}} .
\]
By orthonormality the residual length is
$\norm{\rr_k}^2=\sum_{i=2}^{n} a_i^2\lambda_i^{2k}$. Measure the residual
relative to the parallel part:
\[
\rho_k:=\frac{\norm{\rr_k}}{|a_1|\,|\lambda_1|^{k}}
       =\frac{1}{|a_1|}\sqrt{\sum_{i=2}^{n} a_i^2\Bigl(\tfrac{\lambda_i}{\lambda_1}\Bigr)^{2k}} .
\]
For every $i\ge2$ we have $|\lambda_i/\lambda_1|<1$, so each term
$(\lambda_i/\lambda_1)^{2k}\to0$ and hence $\rho_k\to0$. Thus $\yy_k$ aligns
with $\vv_1$, and after normalisation $\xx_k\to\pm\vv_1$ (the sign tracks
$\operatorname{sign}(a_1\lambda_1^{k})$). \end{proof}

% ---------------------------------------------------------------------
\subsection{(b) Convergence rate and the spectral gap}

\begin{claim}
The residual ratio shrinks each step by at least $|\lambda_2/\lambda_1|$:
$\ \rho_{k+1}/\rho_k\le |\lambda_2/\lambda_1|$.
\end{claim}
\begin{proof}
From the definition of $\rho_k$,
\[
\frac{\rho_{k+1}}{\rho_k}
=\frac{1}{|\lambda_1|}\,
\sqrt{\frac{\sum_{i\ge2} a_i^2\lambda_i^{2k+2}}{\sum_{i\ge2} a_i^2\lambda_i^{2k}}}\, .
\]
Put $\omega_i:=a_i^2\lambda_i^{2k}\ge0$. The fraction under the root is the
weighted average $\bigl(\sum_{i\ge2}\omega_i\lambda_i^{2}\bigr)\big/\bigl(\sum_{i\ge2}\omega_i\bigr)$
of the numbers $\{\lambda_i^2:i\ge2\}$, which cannot exceed their maximum
$\lambda_2^{2}$. Therefore
\[
\frac{\rho_{k+1}}{\rho_k}\le\frac{1}{|\lambda_1|}\sqrt{\lambda_2^{2}}
=\Bigl|\tfrac{\lambda_2}{\lambda_1}\Bigr|<1 .
\]
\end{proof}
Convergence is therefore \textbf{geometric} with rate $|\lambda_2/\lambda_1|$.
The quantity $|\lambda_1|-|\lambda_2|$ (equivalently the ratio
$|\lambda_2/\lambda_1|$) is the \textbf{spectral gap}: a large gap means fast,
robust convergence; a small gap means slow convergence vulnerable to noise.

% ---------------------------------------------------------------------
\subsection{(c) When the spectral gap closes}

\textbf{(i)}~Take $\displaystyle \Xb=\begin{pmatrix}1&0\\0&-1\end{pmatrix}$,
with $\lambda_1=1,\ \lambda_2=-1$, so $|\lambda_1|=|\lambda_2|=1$ (no gap).
For a generic start $\xx_0=(a,b)$ with $a,b\ne0$,
\[
\xx_1=\Xb\xx_0=\begin{pmatrix}a\\-b\end{pmatrix},\quad
\xx_2=\begin{pmatrix}a\\ b\end{pmatrix},\quad
\xx_3=\begin{pmatrix}a\\-b\end{pmatrix},\ \dots
\quad\Longrightarrow\quad
\Xb^{k}\xx_0=\begin{pmatrix}a\\ (-1)^{k}b\end{pmatrix}.
\]
The norm $\sqrt{a^2+b^2}$ is preserved, so after normalisation $\xx_k$
oscillates forever between the two distinct directions
$(a,b)$ and $(a,-b)$ and never converges to a single vector.

\textbf{(ii)}~\extra{explanation not written in the notes.} A vanishing gap is
a \emph{structural} obstruction, not merely slow convergence. The Power Method
works by amplifying each eigencomponent by its own factor $\lambda_i^{k}$; the
top direction wins only because $|\lambda_1|>|\lambda_i|$. When
$|\lambda_1|=|\lambda_2|$ the leading components are amplified by the
\emph{same} factor every step, so their relative weights never change: there is
no selection mechanism. The iterate keeps whatever mixture it started with --
oscillating (as above, when the eigenvalues differ in sign) or, in the
degenerate case $\lambda_1=\lambda_2$, converging to an essentially arbitrary
vector of the eigenspace rather than to a unique eigenvector. In the rate of
part (b), $|\lambda_2/\lambda_1|=1$, so $\rho_k$ does not decay at all.

\textbf{(iii)}~\extra{explanation not written in the notes.} When
$|\lambda_2/\lambda_1|$ is close to (but below) $1$, two problems compound.
First, the number of iterations needed scales like
$1/\bigl(1-|\lambda_2/\lambda_1|\bigr)$, which blows up. Second, each step of
double-precision arithmetic injects rounding error of order $10^{-16}$ into
\emph{all} directions, including $\vv_2$. The useful per-step separation
between $\vv_1$ and $\vv_2$ is only $\propto |\lambda_2/\lambda_1|^{k}$; once that
signal falls to the level of the accumulated noise, the computed vector drifts
within the near-degenerate $\{\vv_1,\vv_2\}$ subspace and the ordering of the
two eigenvalues becomes numerically ambiguous.

% ---------------------------------------------------------------------
\subsection{(d) The two symmetric matrices $\Xb\T\Xb$ and $\Xb\Xb\T$}

Let $\Xb$ be a general (rectangular) $m\times n$ matrix.

\textbf{(i) Symmetric and positive semi-definite.}
\[
(\Xb\T\Xb)\T=\Xb\T(\Xb\T)\T=\Xb\T\Xb,\qquad
(\Xb\Xb\T)\T=\Xb\Xb\T,
\]
so both products are symmetric. For positive semi-definiteness, take any real
vector $\mathbf w$ of the appropriate size:
\[
\mathbf w\T(\Xb\T\Xb)\mathbf w=(\Xb\mathbf w)\T(\Xb\mathbf w)=\norm{\Xb\mathbf w}^2\ge0,
\qquad
\mathbf w\T(\Xb\Xb\T)\mathbf w=\norm{\Xb\T\mathbf w}^2\ge0 .
\]
\fix{the notes wrote $\norm{\Xb\T\mathbf w}^2$ for the $\Xb\T\Xb$ case; the correct
quantity there is $\norm{\Xb\mathbf w}^2$.} Being symmetric (real eigenvalues, by
the spectral theorem) and PSD ($\mathbf w\T A\mathbf w\ge0\Rightarrow$ eigenvalues
$\ge0$), both matrices have \emph{real, non-negative} eigenvalues. Hence
$\sigma_i:=\sqrt{\lambda_i}$ is well defined.

\textbf{(ii) Same nonzero eigenvalues.}
Suppose $\Xb\T\Xb\,\vv=\lambda\vv$ with $\lambda\ne0$ (so $\vv\ne\zero$).
Left-multiplying by $\Xb$,
\[
\Xb(\Xb\T\Xb\,\vv)=\lambda\,\Xb\vv
\quad\Longrightarrow\quad
(\Xb\Xb\T)(\Xb\vv)=\lambda\,(\Xb\vv).
\]
Moreover $\Xb\vv\ne\zero$, since $\Xb\vv=\zero$ would give
$\Xb\T\Xb\vv=\zero=\lambda\vv$, forcing $\lambda=0$. Thus $\Xb\vv$ is an
eigenvector of $\Xb\Xb\T$ with the same eigenvalue $\lambda$. The same argument
with $\Xb\T$ shows every nonzero eigenvalue of $\Xb\Xb\T$ is an eigenvalue of
$\Xb\T\Xb$. Hence $\Xb\T\Xb$ and $\Xb\Xb\T$ have \emph{identical nonzero
eigenvalues}, the shared values being $\sigma_i^{2}$.

% ---------------------------------------------------------------------
\subsection{(e) Power Method for the SVD: algorithm and worked example}

\textbf{(i) Algorithm (language-agnostic pseudocode).} Input: an
$N\times M$ matrix $\Xb$. Output: the top singular triplet
$\sigma_1,\vv_1,\uu_1$.
\begin{enumerate}[leftmargin=2.0em, itemsep=1pt, label=\arabic*.]
  \item \textbf{Form the smaller symmetric matrix.}
        If $N\ge M$, set $A\leftarrow\Xb\T\Xb$ (size $M\times M$);
        otherwise set $A\leftarrow\Xb\Xb\T$ (size $N\times N$).
  \item \textbf{Power iteration on $A$.} Choose a random unit vector $\xx$.
        Repeat $\xx\leftarrow A\xx/\norm{A\xx}$ until
        $\norm{A\xx/\norm{A\xx}-\xx}<\text{tol}$ (or a maximum iteration count).
  \item \textbf{Read off the singular value.}
        $\lambda_1\leftarrow \xx\T A\,\xx$ (Rayleigh quotient);
        $\sigma_1\leftarrow\sqrt{\lambda_1}$.
  \item \textbf{Recover the companion vector.}
        If $A=\Xb\T\Xb$: set $\vv_1\leftarrow\xx\in\R^{M}$ and
        $\uu_1\leftarrow \Xb\vv_1/\sigma_1\in\R^{N}$.
        If $A=\Xb\Xb\T$: set $\uu_1\leftarrow\xx\in\R^{N}$ and
        $\vv_1\leftarrow \Xb\T\uu_1/\sigma_1\in\R^{M}$.
  \item \textbf{Return} $\sigma_1,\vv_1,\uu_1$.
\end{enumerate}
\textit{Choice of matrix.} $A$ is square of size $\min(N,M)$. When $N\ll M$ the
$N\times N$ matrix $\Xb\Xb\T$ is far cheaper to store and to multiply, and when
$M\ll N$ the $M\times M$ matrix $\Xb\T\Xb$ is cheaper -- always pick the smaller.
\fix{the pseudocode text in the notes labelled this rule backwards (``if
$N\gg M$ compute $\Xb\Xb\T$''), although the stated reason and the worked
example below both use the smaller matrix. The consistent rule is given above.}

\textbf{(ii) Worked example by hand.}
\[
\Xb=\begin{pmatrix}1&0\\0&1\\1&1\end{pmatrix}\quad(N=3,\ M=2).
\]
Here $N>M$, so use the $2\times2$ matrix
\[
\Xb\T\Xb=\begin{pmatrix}1&0&1\\0&1&1\end{pmatrix}
\begin{pmatrix}1&0\\0&1\\1&1\end{pmatrix}
=\begin{pmatrix}2&1\\1&2\end{pmatrix}.
\]
\emph{Eigenvalues:}\quad
$\det\!\begin{pmatrix}2-\lambda&1\\1&2-\lambda\end{pmatrix}=(2-\lambda)^2-1
=\lambda^2-4\lambda+3=(\lambda-1)(\lambda-3)=0,$
so $\lambda=1$ or $\lambda=3$; the dominant one is $\lambda_1=3$.

\emph{Eigenvector for $\lambda_1=3$:}\quad $(2-3)x+y=0\Rightarrow y=x$, giving the
unit vector $\displaystyle\vv_1=\tfrac{1}{\sqrt2}\begin{pmatrix}1\\1\end{pmatrix}.$

\emph{Singular value:}\quad $\sigma_1=\sqrt{\lambda_1}=\sqrt3.$

\emph{Left singular vector:}\quad
\[
\uu_1=\frac{\Xb\vv_1}{\sigma_1}
=\frac{1}{\sqrt3}\cdot\frac{1}{\sqrt2}\begin{pmatrix}1&0\\0&1\\1&1\end{pmatrix}\!\begin{pmatrix}1\\1\end{pmatrix}
=\frac{1}{\sqrt6}\begin{pmatrix}1\\1\\2\end{pmatrix}.
\]

\emph{Verification.}
\[
\Xb\vv_1=\tfrac{1}{\sqrt2}\begin{pmatrix}1\\1\\2\end{pmatrix},
\qquad
\sigma_1\uu_1=\sqrt3\cdot\tfrac{1}{\sqrt6}\begin{pmatrix}1\\1\\2\end{pmatrix}
=\tfrac{1}{\sqrt2}\begin{pmatrix}1\\1\\2\end{pmatrix}
\ \Rightarrow\ \Xb\vv_1=\sigma_1\uu_1.\ \checkmark
\]
\[
\Xb\T\uu_1=\tfrac{1}{\sqrt6}\begin{pmatrix}1&0&1\\0&1&1\end{pmatrix}\!\begin{pmatrix}1\\1\\2\end{pmatrix}
=\tfrac{1}{\sqrt6}\begin{pmatrix}3\\3\end{pmatrix}
=\tfrac{3}{\sqrt6}\begin{pmatrix}1\\1\end{pmatrix},
\qquad
\sigma_1\vv_1=\sqrt3\cdot\tfrac{1}{\sqrt2}\begin{pmatrix}1\\1\end{pmatrix}
=\sqrt{\tfrac32}\begin{pmatrix}1\\1\end{pmatrix},
\]
and $\tfrac{3}{\sqrt6}=\sqrt{\tfrac32}$, so $\Xb\T\uu_1=\sigma_1\vv_1.$ \checkmark
Both SVD relations hold, confirming $(\sigma_1,\vv_1,\uu_1)$.

\end{document}
